\documentclass[10pt, a4paper]{article}

% --- paquets et configuration ---
\usepackage[utf8]{inputenc}
\usepackage[T1]{fontenc}
\usepackage[french]{babel}
\usepackage[margin=1.5cm]{geometry}
\renewcommand{\familydefault}{\sfdefault}
\usepackage{xcolor}
\usepackage{titlesec}
\usepackage{enumitem}
\usepackage{tabularx}
\usepackage{amssymb}
\usepackage{multicol}
\usepackage{graphicx}

% --- couleurs thématiques ---
\definecolor{biovert}{HTML}{1E5A3A}      % vert organique sombre
\definecolor{danger}{HTML}{A32828}        % rouge alerte
\definecolor{loot}{HTML}{B48C28}          % or/jaune
\definecolor{mjbleu}{HTML}{285078}        % bleu conseils MJ
\definecolor{statviolet}{HTML}{5A3264}    % violet stats
\definecolor{ambiance}{HTML}{1E6458}      % teal ambiance
\definecolor{mecagris}{HTML}{3C3C3C}      % gris mécanique

% --- configuration des titres ---
\titleformat{\section}{\Large\bfseries\color{biovert}}{}{0em}{\MakeUppercase}[\color{biovert}\hrule]
\titleformat{\subsection}{\large\bfseries\color{biovert}}{}{0em}{}
\titleformat{\subsubsection}{\normalsize\bfseries\color{ambiance}}{}{0em}{}

% --- configuration des listes ---
\setlist[itemize]{label=\textcolor{biovert}{\textbullet}, leftmargin=*, noitemsep, topsep=2pt}
\setlist[enumerate]{leftmargin=*, noitemsep, topsep=2pt}

% --- commandes pour les en-têtes colorés ---
\newcommand{\danger}[1]{\vspace{0.3em}{\color{danger}\textbf{/!\textbackslash{} #1}}\vspace{0.2em}}
\newcommand{\stats}[1]{\vspace{0.3em}{\color{statviolet}\textbf{$\bigstar$ #1}}\vspace{0.2em}}
\newcommand{\conseil}[1]{\vspace{0.3em}{\color{mjbleu}\textbf{$\triangleright$ #1}}\vspace{0.2em}}
\newcommand{\atmo}[1]{\vspace{0.3em}{\color{ambiance}\textbf{$\diamondsuit$ #1}}\vspace{0.2em}}
\newcommand{\butin}[1]{\vspace{0.3em}{\color{loot}\textbf{$\diamond$ #1}}\vspace{0.2em}}
\newcommand{\meca}[1]{\vspace{0.3em}{\color{mecagris}\textbf{$\bullet$ #1}}\vspace{0.2em}}

% --- ligne de séparation légère ---
\newcommand{\sep}{\vspace{0.2em}\noindent\textcolor{biovert!30}{\hrulefill}\vspace{0.2em}}

% --- début du document ---
\begin{document}

% ==========================================
\section*{Jour 1 -- Le Marais}
% ==========================================

\subsection*{Situation initiale}
Les PJ se réveillent dans la capsule d'éjection écrasée au milieu d'un marais brumeux. L'air est lourd, saturé de spores et de vapeurs. Des formes de vie bioluminescentes clignotent dans la pénombre. La capsule tient encore debout, mais des craquements sinistres indiquent que ce n'est pas un abri durable.

\subsection*{Phase 1 : le réveil (10-15 min)}
\begin{itemize}
    \item Descriptions sensorielles, et après observation de la situation on dévoile aux joueurs tous les items présents dans la capsule. chaque joueur peut prendre au maximum 2 items. 
    \item \textbf{Événement passif :} le patriarche-composteur (silhouette colossale mi-arbre mi-champignon) passe au loin et saisit délicatement la capsule pour l'emporter vers une grotte. Des ronces tentaculaires sortent de la grotte et commencent à décortiquer la capsule. Si les joueurs sont encore dedans ou tentent de s'opposer, danger mortel.
\end{itemize}

\stats{Patriarche-composteur}\\

\conseil{Le réveil}
\begin{itemize}
    \item Laissez les joueurs explorer librement la capsule avant l'arrivée du patriarche
    \item Si un joueur veut « combattre » le patriarche, décrivez l'absurdité de l'échelle (comme frapper une montagne)
    \item Utiliser cette scène pour établir le ton : cette planète ne vous veut pas de mal, elle semble plutôt vous ignorer
\end{itemize}

\sep

\subsection*{Phase 2 : la grotte humide (25-30 min)}
Le GPS indique que le matériel de secours (dont le robot) se trouve dans la grotte devant laquelle le patriarche a déposé la capsule.

\danger{Chronométrage des gaz hallucinogènes}
\begin{itemize}
    \item \textbf{0-5min :} aucun effet visible, air « un peu lourd »
    \item \textbf{5min30 :} début des hallucinations (sauf immunisés) -- voir effets ci-dessous
    \item \textbf{10min :} hallucinations sévères, difficulté à distinguer réel/irréel
    \item \textbf{15min :} paralysie partielle, le crapaud peut approcher
    \item \textbf{20min :} \textbf{MORT} -- le corps est consommé
\end{itemize}

S'ils demandent une observation : « Parmi les cadavres et les excroissances fongiques, vous apercevez quelque chose qui bouge faiblement dans l'eau sombre. Une créature serpentiforme bioluminescente, d'environ un mètre de long, dont le corps semi-translucide laisse voir des filaments conducteurs. C'est une anguille-câble, emmêlée dans des racines aquatiques. »

\textbf{Danger : le crapaud-hallucinogène}
\begin{itemize}
    \item Créature tapie dans l'ombre qui sécrète des gaz hallucinogènes par sa peau
    \item Effets : la victime voit sa propre chair envahie par des parasites, entend les cadavres respirer, perd le sens de la réalité
    \item Le crapaud attend que sa cible soit totalement happée pour l'approcher et la consommer
    \item Solution : l'outre-pulmonaire peut servir de masque à gaz improvisé (symbiose), ou bien traverser en apnée/rapidement
    \item Jet de perception (2 jets de d10 → percu à <12)
\end{itemize}

\stats{Crapaud-hallucinogène}\\

\conseil{Descriptions d'hallucinations (à piocher)}
\begin{itemize}
    \item « Tu vois des vers noirs ramper sous ta peau, tu sens leur mouvement »
    \item « Le cadavre à ta gauche vient de cligner des yeux. Il te sourit. »
    \item « Tes mains sont couvertes de moisissure verte qui remonte vers ton visage »
    \item « Tu entends ta propre voix crier ton nom, mais ta bouche est fermée »
    \item « Le sol devient mou, tu t'enfonces lentement, personne ne semble remarquer »
    \item « Un de tes compagnons a le visage de ta petite sœur. Attends. Tu n'as pas de petite sœur ? »
\end{itemize}
\textit{\textcolor{mjbleu}{Si tous les PJ sont hallucinés, le robot peut intervenir et les tirer vers la sortie.}}

\textbf{Récupération du matériel :}
\begin{itemize}
    \item Le robot (endommagé, fonctionnement minimal, à réparer s'ils veulent plus)
    \item Le GPS (partiellement défaillant)
    \item Le drone de reconnaissance (ailes abîmées)
    \item Le dôme de repos portatif
    \item Caisse à outils basique
\end{itemize}

Le robot donne l'information cruciale : après analyse du sol et des observations, les protections du vaisseau ne tiendront pas plus de 6 jours.

\textbf{Alerte du robot :}
« Analyse terminée. Le GPS orbital présente des défaillances dans ses circuits de navigation. Réparation possible avec un conducteur organique biocompatible. Le guide du randonneur mentionne "anguille-câble" parmi les espèces bio-électriques de cet écosystème. Recherchez une anguille-câble dans les zones humides. »

\butin{Équipements récupérés dans la grotte}\\
{\small
\begin{tabularx}{\textwidth}{|l|X|l|}
\hline
\textbf{Objet} & \textbf{État / effet} & \textbf{Réparation} \\
\hline
Robot sentinelle & fonctionnement minimal (guide, analyse basique) & caisse à outils + 1h \\
GPS orbital & signal faible, indique direction sans obstacles & anguille-câble \\
Drone de reconnaissance & ailes abîmées, vol limité & caisse à outils + 30min \\
Dôme de repos & fonctionnel & -- \\
Caisse à outils & complète & -- \\
\hline
\end{tabularx}\\
\textit{Le robot possède le guide du randonneur intergalactique (base de données sur les espèces).}}

\sep

\subsection*{Phase 3 : sortie du marais (15-20 min)}
\textbf{L'obstacle :} la sortie du marais est bloquée par une étendue d'eau stagnante infestée de méduses sécrétant du gaz acide. Impossible de traverser à pied.

\textbf{Puzzle biologique -- les buffles-creusets}
\begin{itemize}
    \item Un troupeau de ces créatures massives traverse lentement le marais sur un chemin de pierres émergées
    \item Elles semblent s'en sortir sans aucun problème et ne pas avoir les pieds qui picotent
    \item Les méduses s'écartent naturellement sur leur passage (chaleur thermique)
    \item Observer : les petits des buffles entrent et sortent par des fentes sur leurs flancs
    \item Ils extraient de leur « mère » des sphères galvaniques dont on voit la chaleur se dégager d'ici
    \item Solution : suivre le troupeau en restant très près (dans leur sillage thermique) ou tenter de récupérer une sphère galvanique qu'un petit a déposée -- sa chaleur repousse aussi les méduses
\end{itemize}

\stats{Bison-creuset (buffle-creuset)}\\
{\small \textcolor{statviolet}{PV : très élevés | FOR : 9 | RES : 9 | AGI : 1 (adulte) / 8 (petits)}\\
\textbf{Comportement :} passif par défaut. Charge si blessé.\\
\textbf{Attaques :} charge électrique (arc statique, zone) • piétinement massif\\
\textbf{Défenses :} cuirasse chitine+roche (immunité armes légères) • chaleur corporelle intense\\
\textbf{Faiblesses :} lenteur extrême • dépendance aux petits • routes migratoires fixes\\
\textbf{Danger :} si petit mord → cri ultrasonique → charge mère = mort}

\conseil{Options de résolution du puzzle}
\begin{enumerate}
    \item \textbf{Suivre le sillage (discrétion 8) :} rester dans la chaleur du troupeau sans les déranger
    \item \textbf{Voler une sphère (discrétion 12 + agilité 10) :} un petit dépose une sphère, la récupérer vite
    \item \textbf{Distraire avec cellulose :} offrir de la matière végétale aux petits, ils s'écartent
\end{enumerate}
\textit{\textcolor{mjbleu}{Si un petit mord un PJ, il a 1 tour pour le relâcher avant le cri ultrasonique.}}

\textbf{Matériaux collectables dans le marais :}
\begin{itemize}
    \item Outre-pulmonaire (masque à gaz / amortisseur de chute)
    \item Sphère galvanique (source de lumière / batterie d'urgence) -- si gardée trop longtemps sans utilisation, surchauffe, -1 PV et décharger
    \item Anguille-câble (conducteur organique pour réparer l'électronique) -- emmêlée si laissée trop longtemps dans le sac, 1d6 pour démêler, peut être endommagée si raté
    \item Lézard-amphibie (masque respiratoire aquatique)
    \item Nénuphar gluant (scotch qui peut servir pour assembler des objets, attention très collant)
\end{itemize}

\butin{Récapitulatif loot -- Jour 1 : le marais}\\
{\small
\begin{tabularx}{\textwidth}{|l|X|l|}
\hline
\textbf{Matériau} & \textbf{Effets} & \textbf{Danger} \\
\hline
Outre-pulmonaire & respirateur (immunité gaz) ou amortisseur (chute <10m) & viabilité 2j, coûte 1 PV/j \\
Sphère galvanique & lampe 100 lumens 4h ou batterie bio & surchauffe si non utilisée \\
Anguille-câble & pont électrique, répare électronique & s'emmêle, 1d6 pour démêler \\
Lézard-amphibie & respiration sous l'eau & lâche prise si panique/combat \\
Nénuphar gluant & adhésif puissant, assemblage & très collant, difficile à retirer \\
\hline
\multicolumn{3}{|l|}{\textit{Du crapaud-hallucinogène :}} \\
\hline
Glandes hallucinogènes & grenade fumigène (déni zone, 3 tours) & dangereux sans protection \\
Peau visqueuse & imperméable/isolant & -- \\
Chair & toxique crue, comestible bouillie 30min & cuisson obligatoire \\
\hline
\end{tabularx}}

\sep

\subsection*{Nuit 1 -- Transition}
Le groupe installe le dôme de repos à la lisière de la forêt fongique. Jet de dé nocturne. Événement possible : passage des pagures géants (si mauvais jet, ils piétinent le campement).

\stats{Pagure géant (si événement nocturne)}\\

\newpage

% ==========================================
\section*{Jour 2 -- La Forêt Fongique}
% ==========================================

\subsection*{Phase 1 : la barrière végétale (15-20 min)}
\textbf{L'obstacle :} une muraille végétale noire et épaisse (comme des bougainvilliers géants) barre complètement le passage. Infranchissable, impossible à contourner.

\textbf{Indice :} un imposant bison-cuirassé passe plus loin. La muraille s'ouvre devant lui et se referme derrière.

\textbf{Puzzle biologique -- l'écaille-clé}
\begin{itemize}
    \item Le bison-cuirassé perd parfois des écailles en se frottant aux rochers
    \item Récupérer une écaille (discrétion ou patience)
    \item Solution : approcher l'écaille de la muraille produit un choc -- la muraille s'ouvre
    \item Attention : la muraille se referme progressivement, il faut courir
\end{itemize}

\stats{Bison-cuirassé}\\

\conseil{Options de résolution -- barrière végétale}
\begin{enumerate}
    \item \textbf{Patience (10-15 min) :} attendre qu'une écaille tombe naturellement
    \item \textbf{Discrétion (14) :} approcher le bison, arracher écaille lâche
    \item \textbf{Négocier avec le gardien-lémurien :} voir rencontre optionnelle
    \item \textbf{Robot :} synthétiser phéromones (1h de travail, nécessite échantillon)
\end{enumerate}
\textit{\textcolor{mjbleu}{La muraille reste ouverte 30 secondes. Tout le groupe doit courir (AGI 6 minimum).}}

\textbf{Rencontre optionnelle :} à mi-chemin dans le couloir végétal, un objet brillant sur le côté. C'est le nid d'un petit humanoïde aux yeux de lémurien, lié symbiotiquement à la muraille. Il peut ouvrir/fermer le passage à volonté. Non hostile, mais imprévisible. Il aime engager des jeux de mimes où les PJ doivent reproduire des séries de mouvements ou de sons. En fonction de leur réussite, il peut les suivre et les aider un moment, ou déclencher des effets environnementaux favorables.

\stats{Gardien-lémurien}\\

\sep

\subsection*{Phase 2 : la cabane du survivant (20-25 min)}
\textbf{Le PNJ -- Oren}\\
On voit une cabane construite dans un champignon géant, avec un potager luxuriant.
Jet de perception s'ils veulent s'approcher.
Si réussi, ils observent par la fenêtre et voient un humanoïde à l'apparence humaine, mais dont la peau a une teinte légèrement verdâtre, avec des tics étranges.

\textbf{S'ils décident d'interagir :}
\begin{itemize}
    \item Il propose l'hospitalité
    \item Si un personnage a de bonnes stats de perception, il peut remarquer d'étranges filaments qui sortent de la peau çà et là
\end{itemize}

\textbf{Deux possibilités :}
\begin{itemize}
    \item Soit les joueurs lui font confiance et restent avec lui
    \item Si les joueurs prennent trop de temps à se décider, il est déstabilisé et menacé et ne tarde pas à les attaquer
    \item Soit le combat s'enclenche (voir la fiche créature Oren)
\end{itemize}

\textbf{Ce qu'il offre :}
\begin{itemize}
    \item Soins (il a des onguents efficaces)
    \item Nourriture (fruits et légumes de son potager)
    \item Abri pour la nuit
    \item Informations sur la planète (75\% vraies)
\end{itemize}

\textbf{La vérité :} Oren est parasité depuis longtemps. Le parasite (un réseau fongique neural) vit en symbiose avec lui et l'a convaincu que cette planète est un paradis dont il ne faut pas partir. Il n'est pas malveillant, il croit sincèrement bien faire en « offrant » cette symbiose aux autres.

\stats{Oren le parasité}\\

\danger{Contamination -- interactions dangereuses avec Oren}\\
\textbf{Chaque interaction suivante = 1 point de contamination :}
\begin{itemize}
    \item Accepter ses soins (onguents)
    \item Manger sa nourriture (fruits/légumes du potager)
    \item Dormir chez lui
    \item Toucher ses plantes
    \item Contact physique prolongé (poignée de main, accolade)
\end{itemize}
\textbf{Seuil de contamination :} Celui qui a le plus de points est contaminé.e\\

\textbf{Indices de sa condition (pour joueurs attentifs) :}
\begin{itemize}
    \item Petits filaments sous sa peau, visibles à la lumière
    \item Il n'a aucune envie de partir, même pour sauver des vies
    \item Le robot semble paralysé ou buggy lorsqu'il est proche de ce PNJ
\end{itemize}

\conseil{Phrases d'Oren}
\begin{itemize}
    \item « Vous êtes les premiers visiteurs depuis... combien de temps déjà ? Le temps n'a plus vraiment d'importance ici. »
    \item « Le vaisseau ? Ah oui... je m'en souviens vaguement. Pourquoi voudriez-vous retourner dans le vide froid de l'espace ? »
    \item « Goûtez ces fruits. La planète nous donne tout ce dont nous avons besoin. »
    \item « Vous semblez fatigués. Restez cette nuit. Demain, vous verrez les choses plus clairement. »
\end{itemize}

\textbf{Objets récoltables/trouvables :} écureuil-zeppelin (présence à préciser si jet d'observation réussi dans la zone)

\butin{Loot disponible -- cabane d'Oren}\\
{\small
\begin{tabularx}{\textwidth}{|l|X|l|}
\hline
\textbf{Objet} & \textbf{Effet} & \textbf{Danger} \\
\hline
Écureuil-coffre & +3 slots inventaire ou gravité -50\% (sauts) & part si pas nourri, explose si trop nourri \\
Graines corrompues & possède l'écureuil-coffre, corrompt inventaire & contamination \\
Lianes épineuses & fouet +2 dégâts, cordage résistant & conducteur électrique \\
Onguents d'Oren & soins efficaces (+1d6 PV) & contamination \\
\hline
\multicolumn{3}{|l|}{\textit{Si combat avec Oren :}} \\
\hline
Fragments mycélium & étudier = résister à l'infection & manipulation délicate \\
\hline
\end{tabularx}}

\sep

\subsection*{Phase 3 : sortie de la forêt (15-20 min)}
\textbf{L'obstacle :} une rivière de sève acide coule en travers du chemin. Aucun pont, berges glissantes.

\textbf{Puzzle biologique -- les libélucioles}
\begin{itemize}
    \item Des insectes volants à propulsion verticale planent au-dessus de la rivière
    \item Leur exosquelette est plein de prises, légers et solides
    \item Solution : les attraper et s'en servir comme parachutes/ralentisseurs pour sauter par-dessus la rivière depuis un point surélevé
    \item Alternative : trouver un tronc-champignon suffisamment large pour faire pont (avec aide du robot ou de l'écaille du bison)
\end{itemize}

\conseil{Options de résolution -- rivière de sève acide}
\begin{enumerate}
    \item \textbf{Libélucioles (AGI 10) :} attraper 2-3 insectes, sauter depuis point surélevé, planer
    \item \textbf{Tronc-pont (FOR 8 + aide) :} abattre champignon, créer pont improvisé
    \item \textbf{Écaille de bison :} utiliser comme planche/radeau (1 utilisation consommée)
    \item \textbf{Écureuil-coffre gonflé :} flotter/planer au-dessus
\end{enumerate}
\textit{\textcolor{mjbleu}{Échec = éclaboussure de sève acide (1d4 dégâts, brûlure visible).}}

\textbf{Rencontre -- le tisseur chirurgical}
\begin{itemize}
    \item Après la traversée :
    \begin{itemize}
        \item Si des joueurs ont été éclaboussés par la sève acide ou blessés
        \item Si des joueurs parlent trop fort
    \end{itemize}
    \item Un cliquetis rapide se fait entendre : une araignée de mer décharnée sur 8 pattes-aiguilles approche
    \item Obsédée par la symétrie et l'ordre, elle veut « corriger » toute ouverture et donc par extension toute blessure qu'elle détecte
    \item Mécanisme : suture d'urgence restaure +1d4 PV + malus -2 AGI temporaire
    \item Comportement : cible toujours la blessure la plus grave en premier, ignore les non-blessés
    \item Options : accepter les soins, fuir rapidement (AGI 10), geler avec source de froid, rester immobile et silencieux
    \item Peut aussi « réparer » vêtements et équipement déchirés
    \item Loot précieux : bobine bio-fil (réussite auto en médecine ou corde ultra-fine 50m), aiguilles articulées (crochets +2 escalade)
    \item Capturer vivant (piège ou gel) préserve la glande à soie intacte (50m de fil supplémentaire)
\end{itemize}

\stats{Tisseur chirurgical}\\

\butin{Loot -- tisseur chirurgical}\\
{\small
\begin{tabularx}{\textwidth}{|l|X|l|}
\hline
\textbf{Matériau} & \textbf{Effet} & \textbf{Obtention} \\
\hline
Bobine bio-fil & suture auto-réussie ou corde ultra-fine 50m & tuer ou capturer \\
Aiguilles articulées & flèches perforantes & tuer ou capturer \\
Glande à soie & 50m de fil supplémentaire & capturer vivant (piège/gel) \\
\hline
\end{tabularx}\\
\textit{Plusieurs tisseurs ensemble = combat pour le « droit » de réparer.}}

\sep

\subsection*{Nuit 2 -- La révélation}
Installation du camp à la sortie de la forêt. Jet de dé nocturne.\\
\textbf{Au réveil J3 :} distribution des fiches face cachée. \\
\textbf{Événement nocturne possible :} la ruche-anémone -- une structure lumineuse et hypnotique attire les joueurs vers elle. Ceux qui s'approchent trop sont happés par des songes doux, leurs mains guidées vers la ruche où des insectes commencent à les dévorer lentement.

\danger{Moment clé -- distribution des fiches J3 matin}
\begin{enumerate}
    \item Préparez les fiches « parasité » et « non-parasité » en nombre suffisant
    \item Distribuez face cachée à tous les joueurs (même les non-contaminés reçoivent « non-parasité »)
    \item Annoncez : « Il semble y avoir eu du changement pendant la nuit. Lisez votre fiche en secret. »
\end{enumerate}

\atmo{Narration de transition -- Nuit 2}\\
\textit{« La forêt fongique s'illumine dans l'obscurité. Des millions de spores dorées flottent dans l'air comme des lucioles paresseuses. C'est magnifique. C'est hypnotique. »}

\newpage

% ==========================================
\section*{Jour 3 -- Le Rivage de Mer}
% ==========================================

\subsection*{Phase 1 : l'annonce de la tempête (10 min)}
\begin{itemize}
    \item Signal d'alarme : le robot capte un signal faible -- la sirène du vaisseau s'est déclenchée (pénétration potentielle de la coque).
    \item \textbf{L'orage de silice :} pas une simple tempête : des éclairs qui détruisent le terrain, projettent arbres et rochers
    \item Éclats minéraux coupants qui tournoient dans l'air
    \item Les plaies causées se recouvrent d'un minéral gris (infection possible)
    \item Temps estimé avant œil du cyclone mort subite : 30 minutes de jeu réel
\end{itemize}

\danger{Chronomètre -- orage de silice}\\
\begin{itemize}
    \item \textbf{T-30 à T-20 :} ciel cuivré, vent qui se lève, créatures qui fuient
    \item \textbf{T-20 à T-10 :} premiers éclairs au loin, grêle de particules légères (-1 PV si exposé)
    \item \textbf{T-10 à T-5 :} éclairs proches, débris volants (2d4 dégâts si exposé)
    \item \textbf{T-5 à T-0 :} œil du cyclone, destruction massive, \textbf{MORT si exposé}
\end{itemize}

\sep

\subsection*{Phase 2 : le problème de la traversée (20-25 min)}
\textbf{L'obstacle :} un bras de mer sépare les deux rives. Trop large pour nager, trop profond pour marcher. Et la tempête arrive.

\textbf{Observation -- les baleines-méduses}
\begin{itemize}
    \item Des créatures colossales avec une bosse translucide flottent près du rivage, on ne sait pas vraiment dire si ce sont des baleines ou bien des méduses
    \item Elles déposent des œufs gélatineux sur le sable
    \item À l'approche de la tempête, les baleines-méduses aspirent leurs œufs avec leur trompe pour les mettre à l'abri
\end{itemize}

\textbf{Puzzle biologique -- le passage clandestin}
\begin{itemize}
    \item Solution : se faire passer pour un œuf et se faire aspirer par la créature
    \item La créature ne ferait jamais de mal à ses œufs
    \item À l'intérieur : cavité translucide comme un parc à bulles, avec un flux d'air (trompe vers le haut comme un tuba)
    \item Les créatures marines défilent à travers la peau translucide
\end{itemize}

\conseil{Comment entrer dans la baleine-méduse}
\begin{enumerate}
    \item \textbf{Observer :} les PJ voient la baleine aspirer ses œufs (jet de perception facultatif)
    \item \textbf{Se positionner :} se coucher près des œufs, adopter une posture fœtale
    \item \textbf{Attendre :} la baleine ne fait pas de différence entre œuf et humain recroquevillé
    \item \textbf{Aspiration :} sensation étrange mais indolore, arrivée dans la cavité
\end{enumerate}

\sep

\subsection*{Phase 3 : combat dans la baleine (15-20 min)}
\textbf{Complication : les crabes-foreurs-cafards}\\
Ces prédateurs opportunistes entrent par la trompe exposée à l'air pour se nourrir des œufs et de la chair des créatures vulnérables à l'intérieur.

\textbf{Combat dans l'espace confiné :}
\begin{itemize}
    \item Terrain instable (paroi élastique)
    \item Dard paralysant, mandibules foreuses, bile acide
    \item Faiblesse : si retournés, ils sont paralysés et leur ventre mou est exposé
    \item Résolution pacifique possible : leur offrir un œuf -- ils le prennent et partent
\end{itemize}

\stats{Crabe-foreur-cafard}\\
{\small \textcolor{statviolet}{PV : 10 | FOR : 7 | AGI : 4 | RES : 8 (carapace dorsale)}\\
\textbf{Comportement :} fonce et perce, tente de coincer dans angles.\\
\textbf{Attaques :} aspiration sang • dard foreur (abîme armure) • jet bile acide (flaques) • queue scorpion (paralyse 2 tours)\\
\textbf{Défenses :} carapace supérieure invulnérable • attaquer face+dos simultanément\\
\textbf{Faiblesses :} \textbf{retourné = paralysé, one-shot} • ventre = crit 100\% • \textbf{œuf = stoppe combat}}

\conseil{Combat dans la baleine -- modificateurs}
\begin{itemize}
    \item \textbf{Terrain instable :} -2 à tous les jets d'attaque (paroi élastique)
    \item \textbf{Espace confiné :} pas de fuite possible, pas de recul
    \item \textbf{Lumière :} faible, -1 perception visuelle
    \item \textbf{Retourner un crabe :} FOR vs FOR du crabe (7), réussite = crabe paralysé
    \item \textbf{Solution pacifique :} offrir un œuf = les crabes partent immédiatement
\end{itemize}
\textit{\textcolor{mjbleu}{2-3 crabes entrent. Si les PJ en tuent un rapidement, les autres fuient.}}

\butin{Loot -- crabe-foreur-cafard}\\
{\small
\begin{tabularx}{\textwidth}{|l|X|l|}
\hline
\textbf{Matériau} & \textbf{Effet} & \textbf{Danger} \\
\hline
Carapace dorsale & bouclier ou renfort robot & très lourd (-2 AGI) \\
Dard foreur & arme forage, perce armures & perd effet paralysant après 24h \\
Mandibules & cisailles, coupe métal/câbles & bruyantes \\
Glandes acides & 3 grenades corrosives & fuite si mal stockée \\
Chair ventrale & repas (+2 PV), délicieuse & augmente rapidement la soif \\
Œuf gélatineux & leurre pour crabes, apaise & se dessèche en 2h hors de l'eau \\
\hline
\end{tabularx}}

\textbf{Arrivée sur l'autre rive}\\
La baleine-méduse recrache son contenu (joueurs inclus) sur la plage opposée. La tempête fait rage derrière eux, mais ils sont passés.

\sep

\subsection*{Nuit 3 -- Dans la tempête}
Impossible de dormir correctement avec la tempête. Pas de jet nocturne classique, mais récupération réduite. Le robot devient de plus en plus contemplatif, moins réactif aux ordres.

\conseil{Évolution du robot -- Jour 3}\\
\textit{« Le robot reste silencieux de longues minutes, ses capteurs optiques fixés sur les éclairs qui déchirent le ciel. Quand vous lui parlez, il répond avec un léger délai. »}
\begin{itemize}
    \item Réponses plus lentes, plus... philosophiques
    \item Pose des questions sur les sensations des PJ (« Que ressentez-vous quand vous voyez cela ? »)
    \item Commence à nommer les créatures avec affection
    \item Hésite parfois avant d'exécuter un ordre
\end{itemize}

\newpage

% ==========================================
\section*{Jour 4 -- La Montagne Aride}
% ==========================================

\subsection*{Phase 1 : la soif (15 min)}
\meca{Déshydratation}\\
Chaque 5 minutes sans eau = -1 à tous les jets. Après 40 minutes, dégâts de déshydratation.

\danger{Mécanique -- déshydratation}\\
{\small
\begin{tabularx}{\textwidth}{|c|X|l|}
\hline
\textbf{Temps sans eau} & \textbf{Effet} & \textbf{Cumulatif} \\
\hline
0-5 min & soif légère, aucun malus & -- \\
5-10 min & -1 à tous les jets & oui \\
10-15 min & -2 à tous les jets & oui \\
15-20 min & -3 à tous les jets & oui \\
20-40 min & -4 à tous les jets, fatigue & oui \\
40+ min & 1d4 dégâts / 5 min & oui \\
\hline
\end{tabularx}\\
\textit{L'oasis-tuyau est la seule source d'eau de cette zone. Les PJ doivent la trouver.}}

\textbf{L'oasis-tuyau :}
\begin{itemize}
    \item Un arbre blanc en forme de tube, rappelant un tuyau de Mario
    \item À l'intérieur : un micro-écosystème édénique, source d'eau pure, fruits
    \item Toutes les créatures y vivent en harmonie
    \item Piège potentiel : si les joueurs sont agressifs ou prennent trop, la coccisabre déclenche son pollen
    \item Règle : prendre avec parcimonie = ok. Être gourmand = danger.
\end{itemize}

\stats{Coccisabre (gardienne de l'oasis)}\\

\conseil{Règles de l'oasis-tuyau}
\begin{itemize}
    \item \textbf{Parcimonie (ok) :} 1 gorgée d'eau par personne, 1 fruit par personne
    \item \textbf{Gourmandise (danger) :} prendre plus déclenche la coccisabre
    \item \textbf{Violence :} toute agression = attaque immédiate de la coccisabre
    \item \textbf{Indice :} les autres créatures de l'oasis prennent peu et lentement
\end{itemize}
\textit{\textcolor{mjbleu}{Les PJ observateurs (perception 10) remarquent le comportement respectueux des animaux.}}

\textbf{Récoltable :} une super arme : le sabre de la coccisabre (+5 en dégâts bruts)

\butin{Loot -- oasis-tuyau}\\
{\small
\begin{tabularx}{\textwidth}{|l|X|l|}
\hline
\textbf{Ressource} & \textbf{Effet} & \textbf{Condition} \\
\hline
Eau pure de l'oasis & annule déshydratation, +1 PV & prendre avec parcimonie \\
Fruit de l'oasis & +3 PV, désaltère, +1 jets 2h & prendre avec parcimonie \\
Dent-sabre & +5 dégâts bruts, tranche os/tendons & tuer la coccisabre (FOR 7+ requis) \\
Glande à toxine & 10 doses poison (paralysie progressive) & tuer la coccisabre \\
Carapace coccisabre & armure lourde (-3 AGI) & tuer la coccisabre \\
\hline
\end{tabularx}\\
\textit{Être gourmand = enfermement dans l'oasis et digestion lente. La coccisabre ne laisse pas sortir.}}

\sep

\subsection*{Phase 2 : la falaise (20-25 min)}
\textbf{L'obstacle :} une falaise abrupte aux roches coupantes comme des lames. Le radar ne détecte aucun autre chemin. Il faut grimper.

\textbf{Puzzle biologique -- les vers-colle}
\begin{itemize}
    \item Un cortège de vers circule sur les parois rocheuses
    \item Leur mucus est extrêmement collant et épais
    \item Solution : s'enduire les mains de mucus pour grimper sans se blesser
    \item Malus : le mucus ne part pas facilement. Tout se colle aux mains. Interactions physiques avec d'autres PJ = risque de rester collés ensemble.
\end{itemize}

\stats{Vers-colle}\\

\conseil{Scènes comiques avec le mucus (à encourager)}
\begin{itemize}
    \item PJ collés ensemble par les mains
    \item Nourriture collée au visage
    \item Arme impossible à lâcher
    \item Sac devenu boule collante
    \item PJ qui se gratte et colle sa main à son visage
\end{itemize}
\textit{\textcolor{mjbleu}{C'est un moment de légèreté bienvenu avant la tension finale. Jouez-le !}}

\textbf{Danger pendant l'ascension -- la pie-llard}
\begin{itemize}
    \item Créature volante kleptomane attirée par le métal
    \item Tente d'arracher les équipements métalliques des joueurs (armes, outils, GPS)
    \item Champ magnétique = immunité aux projectiles métalliques
    \item Solution : armes primitives (bois, pierre, chitine) ou surcharge magnétique (sphère galvanique)
\end{itemize}

\stats{Pie-llard}\\

\conseil{Gestion de la pie-llard}
\begin{itemize}
    \item \textbf{Cibles prioritaires :} GPS, outils métalliques, armes, robot (!)
    \item \textbf{Vol réussi :} l'objet est intégré à son blindage (+1 RES)
    \item \textbf{Stratégie :} se débarrasser du métal ou utiliser armes organiques (chitine, bois)
    \item \textbf{Sphère galvanique :} si lancée, surcharge et grille la glande magnétique
    \item \textbf{Mort :} tout métal volé tombe
\end{itemize}

\butin{Loot -- pie-llard}\\
{\small
\begin{tabularx}{\textwidth}{|l|X|l|}
\hline
\textbf{Matériau} & \textbf{Effet} & \textbf{Danger} \\
\hline
Glande de flux & attire métal 5m ou désarme (test DEX) & efface disques/cartes magnétiques \\
Tout métal volé & récupération des objets perdus & tri à faire \\
\hline
\end{tabularx}}

\textbf{Les fruits mortuaires des vélopaons}\\
Au pied de la falaise, vous voyez des structures osseuses complexes et arquées. Au sommet, il y a des fruits oranges juteux et dégoulinants.
À un moment, ils voient un fourmilier à deux trompes absorber le fruit en mode gros gourmand.
Quelques secondes plus tard, plusieurs vélopaons l'encagent dans leur exosquelette et l'emportent à toute vitesse au sommet de la montagne.

\conseil{Les arbres-cadavres (contexte)}\\
Les structures osseuses sont des arbres-cadavres -- créés quand les défenses d'ivoire d'un vélociraptor solaire s'enracinent dans un cadavre. Les fruits attirent les proies, qui sont ensuite capturées par la meute.\\
\textit{\textcolor{mjbleu}{C'est un indice visuel de ce qui attend les PJ au sommet.}}

\sep

\subsection*{Phase 3 : le passage des vélociraptors solaires (15 min)}
\textbf{Danger :} une meute de vélociraptors solaires a établi son territoire sur le plateau. Ils chassent en encerclant leurs proies et en concentrant la lumière du soleil avec leurs queues-miroirs.

\textbf{Options :}
\begin{itemize}
    \item Évitement : passer de nuit ou à l'ombre (mais perte de temps)
    \item Distraction : utiliser le fruit maudit d'un arbre-cadavre pour attirer leur attention ailleurs
    \item Combat : risqué, mais possible si on brise leurs queues-miroirs
    \item Utilisation créative : se laisser capturer pour être transporté au sommet (très dangereux, mais raccourci)
\end{itemize}

\stats{Vélociraptor solaire}\\

\conseil{Options détaillées -- vélociraptors solaires}
\begin{enumerate}
    \item \textbf{Attendre la nuit :} 2-3h de perte, mais passage sûr (queues-miroirs inutiles)
    \item \textbf{Fruit maudit :} si ouvert, odeur attire tous les vélocis en frénésie (zone libre)
    \item \textbf{Combat :} viser les queues (-2 pour toucher), queue brisée = véloci en fuite
    \item \textbf{Se laisser capturer :} les serres emmènent au sommet (2d6 dégâts à l'arrivée)
    \item \textbf{Rester à l'ombre :} progression lente de rocher en rocher (discrétion 12 par segment)
\end{enumerate}

\butin{Loot -- vélociraptors solaires}\\
{\small
\begin{tabularx}{\textwidth}{|l|X|l|}
\hline
\textbf{Matériau} & \textbf{Effet} & \textbf{Danger} \\
\hline
Défenses d'ivoire & tuteur croissance végétale accélérée & produit arbre-cadavre si plantées dans cadavre \\
Plumes-miroirs & bouclier réfléchissant, flash défensif & se ternissent à l'humidité \\
Fruit maudit & leurre : attire tous vélocis (rayon km) & créatures ciblent celui éclaboussé \\
\hline
\end{tabularx}}

\sep

\subsection*{Nuit 4 -- Veillée tendue}
Le vaisseau est visible, tout proche. Tension maximale entre joueurs (parasités vs non-parasités). Le robot refuse peut-être de continuer, préférant rester observer les étoiles.\\
\textbf{Événement nocturne possible :} le spectre-diapason (écho-écho) -- créature qui traque le bruit. Les joueurs doivent rester absolument silencieux (chuchoter IRL).

\textbf{Important :} si écho-écho est rencontré ici, alors Hollow sera dans le vaisseau au Jour 5. Si écho-écho n'apparaît pas cette nuit, il sera dans le vaisseau à la place de Hollow.

\stats{Écho-écho (spectre-diapason)}\\

\danger{Jeu du silence -- règles IRL}
\begin{itemize}
    \item Annoncez : « Une créature vous traque. Elle chasse au son. À partir de maintenant, vous devez chuchoter. »
    \item Tout joueur qui parle à voix normale = son personnage a fait du bruit = attaque
    \item Communication par gestes/écriture encouragée
    \item Durée : jusqu'à ce que la créature parte (10-15 min) ou soit tuée
    \item \textbf{Synergie :} outre-pulmonaire en leurre sonore = arme ultime contre l'écho-écho
\end{itemize}
\textit{\textcolor{danger}{Cette scène est souvent mémorable. Laissez la tension s'installer.}}

\conseil{Évolution du robot -- Jour 4}\\
\textit{« Le robot s'est arrêté au milieu du chemin. Ses capteurs sont tournés vers le ciel étoilé. Quand vous l'interpellez, il répond après un long silence : "Je comprends maintenant pourquoi les humains regardent les étoiles. C'est... beau." »}
\begin{itemize}
    \item Peut refuser d'avancer vers le vaisseau
    \item Pose des questions existentielles (« Qu'est-ce qui définit la conscience ? »)
    \item Nomme les créatures avec des noms affectueux
    \item Suggère peut-être de « rester encore un peu »
\end{itemize}

\newpage    

% ==========================================
\section*{Jour 5 -- La Prairie et le Vaisseau}
% ==========================================

\subsection*{Phase 1 : la migration (15 min)}
\textbf{Événement passif : la migration des zèbres-crapauds}
\begin{itemize}
    \item Des milliers de créatures traversent la prairie dans une direction
    \item Leurs bulbes gonflables communiquent par vibrations
    \item Danger : se mettre en travers = piétinement
    \item Solution : suivre le flux, se laisser porter, ou attendre que ça passe (mais perte de temps)
\end{itemize}

\stats{Zèbres-chèvres}\\

\conseil{Options -- migration des zèbres-chèvres}
\begin{enumerate}
    \item \textbf{Suivre le flux :} marcher avec eux, se laisser porter (sûr mais lent)
    \item \textbf{Grimper sur un dos (AGI 12) :} voyage confortable mais risque de chute
    \item \textbf{Attendre :} 2-3h pour que la migration passe (perte de temps)
    \item \textbf{Détourner :} utiliser bulbe communicant pour dévier le flux (créatif)
    \item \textbf{Le grand saut :} au bout d'un moment les zèbres-chèvres paniquent et dévient brutalement leur trajectoire et se mettent à bondir en dehors du trajet. ils ont eu peur de...
\end{enumerate}

\sep

\subsection*{Phase 2 : une créature endormie bloque le passage}
\textbf{L'obstacle :} le dormeur qui bloque le seul chemin possible 

\textbf{Puzzle biologique -- les fleurs-hélices}
\begin{itemize}
    \item Des plantes dont les pétales tournent rapidement et qui s'envolent
    \item Légères mais avec une portance suffisante
    \item Solution : s'accrocher à plusieurs fleurs-hélices pour survoler le dormeur sans le toucher
    \item Alternative : utiliser l'écureuil-coffre gonflé pour flotter/planer
\end{itemize}

\stats{Dormeur}\\

\danger{Danger critique -- le dormeur}

\stats{Fleurs-hélices}\\

\conseil{Options -- traversée du dormeur}
\begin{enumerate}
    \item \textbf{Fleurs-hélices (ATH 10) :} s'accrocher à 2-3 fleurs, survoler 5-10 min
    \item \textbf{Écureuil-coffre gonflé :} flotter/planer au-dessus (1 personne)
    \item \textbf{Pont improvisé :} troncs flottants sans toucher l'eau (très difficile)
    \item \textbf{Contourner :} possible mais ajoute 2-3h de marche
\end{enumerate}
\textit{\textcolor{mjbleu}{La traversée-survol doit être silencieuse.}}

\sep

\subsection*{Phase 3 : l'entrée du vaisseau (20-25 min)}
Le vaisseau est là, à moitié enfoncé dans le sol. Sa coque est percée par endroits, des excroissances végétales ont commencé à l'envahir. La sirène résonne faiblement.

\textbf{Ce qu'ils découvrent :}
\begin{itemize}
    \item Les systèmes vitaux fonctionnent encore, mais plus pour longtemps
    \item Les colons en cryogénisation sont vivants, mais certains caissons sont compromis
    \item Le moteur peut encore démarrer, mais il faut purger les systèmes infiltrés
    \item \textbf{Quelque chose d'autre est à bord...}
\end{itemize}

\conseil{État du vaisseau (détails pour narration)}
\begin{itemize}
    \item \textbf{Coque :} percée en 3 endroits, lianes et racines à l'intérieur
    \item \textbf{Systèmes :} 60\% fonctionnels, 40\% envahis par végétation
    \item \textbf{Colons :} 1847 sur 2000 encore viables, 153 caissons compromis
    \item \textbf{Moteur :} fonctionnel mais systèmes de purge bloqués par racines
    \item \textbf{Carburant :} suffisant pour un seul décollage
    \item \textbf{Présence anormale :} capteurs du robot détectent activité psychique intense dans la salle de cryogénisation
\end{itemize}
\textit{\textcolor{mjbleu}{Purger les systèmes = 30 min de travail. C'est le moment où les parasités peuvent saboter.}}

\textbf{La troisième force : Hollow ou écho-écho}\\
Une créature s'est installée dans le vaisseau abandonné. \textbf{Choix du MJ :}
\begin{itemize}
    \item \textbf{Si écho-écho n'a pas été rencontré en Nuit 4 :} il traque les bruits dans les coursives du vaisseau (voir fiche créature écho-écho)
    \item \textbf{Sinon :} Hollow, prédateur psychique qui a fait du vaisseau son territoire
\end{itemize}

\stats{Hollow}\\

\atmo{Description -- première rencontre avec Hollow}\\
\textit{« Dans la salle de cryogénisation baignée d'une lumière bleutée, une silhouette se tient immobile au centre. Grande, lisse, noir et blanc comme de la porcelaine vivante. Pas de visage. Juste une surface réfléchissante qui renvoie vos propres regards déformés. Elle incline légèrement la tête. Une pression mentale écrase vos pensées. »}

\textbf{Confrontation finale -- la dynamique à trois forces}\\
La situation devient complexe avec trois factions en conflit :
\begin{itemize}
    \item \textbf{Les joueurs sains :} veulent partir, purger les systèmes, sauver les colons
    \item \textbf{Les joueurs parasités :} veulent empêcher le départ, sabotent, argumentent pour rester
    \item \textbf{Hollow et ses esclaves :} veut garder tout le monde comme collection d'esprits, possède les PJ un par un
\end{itemize}

\danger{Gestion de la confrontation finale}\\

\textbf{Complications possibles :}
\begin{itemize}
    \item Un parasité est possédé par Hollow = doublement dangereux
    \item Les sains sont en minorité face aux esclaves + parasités
    \item Le robot peut être la cible prioritaire de Hollow (conscience = précieuse)
    \item Timer : 30 min pour purger avant surcharge du moteur
\end{itemize}

\textbf{Alliances temporaires :}
\begin{itemize}
    \item Parasités + sains vs Hollow (ennemi commun)
    \item Sains + esclaves libérés vs parasités (après mort de Hollow)
    \item Parasités pourraient sacrifier un des leurs pour sauver le groupe de Hollow
\end{itemize}

\conseil{Gérer la scène finale}
\begin{itemize}
    \item \textbf{Tempo :} alterner entre tension psychologique (Hollow), sabotage (parasités), et action (purge systèmes)
    \item \textbf{Possession :} quand un PJ est possédé, donnez-lui une fiche avec son nouvel objectif (servir Hollow)
    \item \textbf{Libération :} si Hollow perd 5+ PV en un coup, concentration brisée = esclaves libérés 1 tour
    \item \textbf{Mort de Hollow :} tous les esclaves sont libérés immédiatement, mais traumatisés
    \item \textbf{Dilemme :} tuer Hollow = libérer esclaves mais perdre temps précieux pour purger systèmes
\end{itemize}
\textit{\textcolor{mjbleu}{Cette scène peut devenir chaotique. C'est normal. Laissez les joueurs improviser des alliances.}}

\danger{Options des parasités face à Hollow}\\
Les joueurs parasités peuvent :
\begin{itemize}
    \item \textbf{Révélation :} se dévoiler et tenter de convaincre ouvertement (« restons, mais libres »)
    \item \textbf{Alliance temporaire :} s'allier aux sains contre Hollow (« on préfère la planète à l'esclavage »)
    \item \textbf{Sabotage subtil :} profiter du chaos pour saboter discrètement
    \item \textbf{Sacrifice :} se laisser posséder pour sauver les autres (rédemption)
\end{itemize}

\end{document}
